\documentclass[]{elementary-physics}

\title{5. Increasing Energy}
\date{2023-01-05 v0.1}

\begin{document}

\maketitle

\tableofcontents

\section{Abstract}

[TODO rename file]

\section{Background}

In physics and chemistry, the law of conservation of energy\cite{energy} states that the total energy of an isolated system remains constant; it is said to be conserved over time.

This implies that the sum of all powers in the system equals to zero which is guaranteed by Kirchhoff's circuit laws\cite{Kirchhoff} and Newton's 3rd law of motion\cite{Newton}.

But what about energy transfer? Is the amount of energy leaving one system always the same as the amount of energy entering another? Would double-entry\cite{DoubleEntry} bookkeeping on power always balance?

For instance; consider the equation of a harmonic oscillator\cite{oscillator}, a LRC-circuit driven by an AC-source expressed as a sum of voltages

\begin{equation}
Lq'' + Rq' + C^{-1}q - sin(\omega t) = 0
\end{equation}

Once solved we can change the equation from a sum of voltages to a sum of powers by multiplying with $q'$

\begin{equation}
Lq''q' + Rq'q' + C^{-1}qq' - sin(\omega t)q' = 0
\end{equation}

But two equations are missing; one that contains $-Rq'q'$ where the heat is going to and one that contains $sin(wt)q'$ where the power running the generator is coming from, now we have a system of equations:


\begin{subequations}
\begin{align}
other.terms - Rq'q' = 0 \\
Lq''q' + Rq'q' + C^{-1}qq' - sin(\omega t)q' = 0 \\
other.terms + sin(\omega t)q' = 0
\end{align}
\end{subequations}


We will ignore the heat and focus on the electro-mechanicals; a wheel with magnets and stationary coils. The two equations involved are Faraday’s Law\cite{Faraday} $V = Blv$ on the electrical side and Lorentz force\cite{Lorentz} $F = BlI$ on the mechanical side. They eace replace one of the two $sin(\omega t)$ where B is now a function $B(\theta)$.

We get

\begin{subequations}
\begin{align}
Lq'' + Rq' + C^{-1}q - Bl\theta' = 0 \\
J\theta'' + r\theta' + Blq' = 0
\end{align}
\end{subequations}

where $J$ is the inertia of the wheel and $r$ is the friction. If we change the equations to powers

\begin{subequations}
\begin{align}
Lq''q' + Rq'q' + C^{-1}qq' - Bl\theta'q' = 0 \\
J\theta''\theta' + r\theta'\theta' + Blq'\theta' = 0
\end{align}
\end{subequations}

we have a system where the two equations communicate energy in a balanced way. In this paper we question that it is always so, experiments suggests that the two $B$-functions are usually but not always the same.

\section{Simulations}

To aid us in our research we have used a free simulator called FEMM\cite{FEMM} which can be scripted using Lua 4.0\cite{Lua}. There is also a command-line version of FEMM called xfemm\cite{xfemm} that we have used heavily.

The femm programs allows us to create snapshots of set-ups with coils and magnets. The simulation then provides us with torque and flux-linkage.

[TODO check if still valid and rewrite for latex]

 1: select a wire-diameter D
 2: calculate 
    number of turns N = 56\% * coil-area / wire-area
    resistance R = pCu * 1000 * wire-length / wire-area 
      where pCu = 1.72E-8
    currrent for 1 W, I = sqrt(1 / R)
 3: test the coil twice (with and without current) for each angle
 4: get impulse T = the sum of all torque (with current minus without current) times delta-angle (integrate)
 5: get the flux-linkage when no current
 6: get the flux-linkage divided by delta-angle to get the voltage (derivate)
 7: get impulse V = the sum of all voltage times delta-angle (integrate)
 8: assuming linearity we calculate for each coil
    Tk = T/N*I
    Vk = V/N
    Rk = R/N\^2
    Nk = 56\% * coil-area * 4/pi
    VpT = Vk / Tk
 9: VpT > 1 is a generator-coil and VpT < 1 is a motor-coil
10: we can now choose any wire diameter D and calculate
    N = Nk/D\^2
    T = Tk*N*I
    V = Vk*N
    R = Rk*N\^2
11: for each pair of coils, mo and gen, with any wire, moD and genD, calculate a relative torque for comparisom between coil-pairs
    T = moT - genT = (moTk*moN - genTk*genN)*I
      where I = (genV - moV) / (genR + moR)



\section{Algorithms}

[TODO Check code and algorithms in Lua and FEMM]
[Maybe by Odd Larsson?]

\section{Experiments}

[TODO Validate the simulations with experiments]

\section{Conclusion}

[TODO]

\appendix

\section{In Plain English}

[TODO]

\section{På Ren Svenska}

[TODO]

\section{This Paper}

This is one paper from a collection of papers, all free to be downloaded and shared. If you have ideas how to enhance any of the papers, if you want to contribute, don’t hesitate to contact us at \url{hob.nilre@gmail.com}.\\
\\
The papers can all be found at:
\begin{itemize}
\item \url{https://sites.google.com/site/nilrehob/home/elementary-physics}
\item \url{https://independent.academia.edu/HobNilre/Papers}
\item \url{https://groups.yahoo.com/neo/groups/EVGRAY/files/Hob/}
\item \url{http://overunity.com/15796/elementary-physics-revisited/}
\item \url{http://idipsum.se/home/elementary%20physics.html}
\item \url{https://github.com/boherlin/elementary-physics/tree/master/pdf}
\end{itemize}

They are updated with new versions in an unpredictable manner, possibly not on all sites but at least on the last two sites in the list, make sure you always have the latest version!
Their \LaTeX source-codes can be found at \url{https://github.com/boherlin/elementary-physics/tree/master/src}.
All papers, but not all versions, have been stamped at \url{http://www.OriginStamp.org}.\\
\\
If you enjoyed this paper, found value in it and want to help us, please consider giving us a donation in bitcoin, this is our address:

\begin{figure}[ht] \centering
	\includegraphics[]{1B79p75vQw4Rb1GQdmGYpDapFwEytFJDqw} \caption{1B79p75vQw4Rb1GQdmGYpDapFwEytFJDqw}
\end{figure}


[TODO]

\printbibliography

\end{document}
